% Using InitechBase, Appendix C: Structure of InitechBase Files (annex Sec
% 8.2 UC; PUB-REF-0001 CP-11: the .DBF/.DBT structure from
% spec/samir/dbf_format.h).  Every offset and size is generated from the
% locked headers by harness/pubs/samir_tables.c formats (samir_formats);
% the .FRM description follows os/samir/cmd/report.c rp_load and
% harness/pubs/samir_sample_disk.c.
\chapter{Structure of InitechBase Files}
\Gates{test-dbf-header, test-dbf-fields, test-dbt-read, test-ndx-parse, test-samir-report}
\NoDictionaryHeads
\TableRows{samir_formats}
This appendix describes the layout of the database, memo and index files, for
programmers who read or write them with other programs.  Numbers are stored
with the least significant byte first.  Offsets are hexadecimal.

\section{Database Files (.DBF)}
A database file has three parts: the header, one field descriptor for each
field, and the records.

\ATEHead{Header}
\FormatTable{dbfhdr}

\ATEHead{Field Descriptor}
Each field is described by \FormatDBFRecordBytes{} bytes; the descriptors
follow the header in the order of the fields, and the byte \FormatDBFTerminator{}
follows the last.
\FormatTable{dbfdesc}

\ATEHead{Records}
The records follow the header and the descriptors without a gap, each of the
length the header gives.  The first byte of a record is its deletion mark:
\FormatDBFLive{} (a blank) for an ordinary record, \FormatDBFDeleted{} (an
asterisk) for a record marked for deletion.  The fields follow in the order
of the descriptors, each as ASCII characters of its width: character fields
padded with blanks on the right; numeric fields right-aligned with leading
blanks; dates as eight digits, year, month and day (\texttt{YYYYMMDD}), or eight
blanks; logical fields as \texttt{T} or \texttt{F}; memo fields as the number
of the first block of the text in the memo file, right-aligned in ten
characters, or ten blanks when the record has no memo.

\section{Memo Files (.DBT)}
A memo file is divided into blocks of \FormatDBTBlock{} bytes.  The first block
holds, in its first four bytes, the number of the next free block.  The text
of a memo begins at the start of a block and ends with the two bytes
\FormatDBTTerm{}; a text longer than one block continues in the blocks that
follow.  Lines of the text are separated by a carriage return and a line feed.

\section{Index Files (.NDX)}
An index file is divided into pages of \FormatNDXPage{} bytes.  The first page
is the header:
\FormatTable{ndxhdr}

The other pages hold the keys, in order, each with the record number it
belongs to.  A character key is the value of the key expression, padded with
blanks to the key length; a numeric or date key is the value as an eight-byte
floating-point number.

\section{Report Form Files (.FRM)}
A report form file is 1990 bytes long.  It holds the title of the report, the
group and subgroup expressions with their headings, and for each column its
width, its decimal places, whether it is totalled, its expression and its
heading, followed by the page width, the lines on a page and the margins.
\cmd{REPORT FORM} refuses a file of any other form with
\xmsg{50}{Report file invalid.}
